%======================================================================
\section{Summary}\label{sec:summary}
%======================================================================

\textbf{Goal.} BruceWorx (FreeCAD) reaches \emph{full parity with Fusion~360 and SolidWorks} across the
whole product lifecycle, so that it can become the CAD system of record for BruceWorx products. That covers
sketching and part modeling, assembly, simulation, ASME drawings, part and drawing management, and the
files sent to the 3D printer, the CNC shop and the laser cutter. This replaces the earlier position
(FC-T01) that FreeCAD only complements Fusion. The switch itself is the last milestone (M6) and is gated on the parity
matrix in Appendix~\ref{app:parity}.

\textbf{How this document was built.} The previous version of this file was a UI plan (Fusion-style
layout, patterns, command line). It is now one chapter's worth of material inside a roadmap with
\textbf{13 tracks and 87 epics}. The inputs were:
\begin{itemize}
  \item a five-round interview with you (18--19 Oct 2026), recorded in \S\ref{sec:decisions};
  \item the BruceWorx Ops ``work'' list (191 items, marked \emph{sample data}), used as \emph{ideas only}.
        Appendix~\ref{app:ideas} maps every idea group to the epic that absorbs it;
  \item a capability survey of FreeCAD~1.1.4 in this repository (grep of the actual sources and headless
        experiments), compared with Fusion~360 and SolidWorks~2026 public material;
  \item the repository's merged work (\code{[BW]} commits \#3--\#9), which marks four epics done and five in
        progress.
\end{itemize}

\textbf{Granularity.} The roadmap stops at \emph{epics}. Each epic has a scope, dependencies, an acceptance test and
the ideas it absorbs. Tasks are broken down when an epic starts. \file{roadmap-epics.csv} carries the same
87 epics for import into Ops. It is generated from \file{scripts/epics.py}, together with the tables in this
document, so the two cannot drift.

\textbf{Styling comes later.} The UI and styling track runs after M2. The figures here are working renders
and will be regenerated for the new UI (\eref{UIX-09}).

\subsection{Decisions from the interview}\label{sec:decisions}
\begin{center}
\small
\begin{longtable}{@{}p{3.6cm}p{12.4cm}@{}}
\toprule
Topic & Decision\\
\midrule\endhead
End state & Full Fusion~360 / SolidWorks parity; FreeCAD becomes system of record at M6\\
Ops task list & Ideas only; fresh IDs (\code{BW-\textit{TRK}-\textit{nn}}), origins kept in Appendix~\ref{app:ideas}\\
Manufacturing & FDM printing (PETG/ABS), CNC machining, sheet metal / laser. Not injection molding\\
First stage & Part and assembly modeling (main track)\\
Track order & Modeling/assembly first; \textbf{PDM evaluation and drawings run in parallel from the start}\\
Core modeling & Configurations; direct edit on imported STEP; parameters and top-down; threads and hardware\\
Advanced modeling & Freeform / T-spline-like; surfacing parity; mesh $\rightarrow$ solid; frames / weldments\\
File model & SolidWorks-style: \textbf{one file per part}, each with its own part number; assemblies link to files\\
Assembly extras & In-context editing; large-assembly mode; smart fastener mating; dynamic motion study\\
Simulation & Static + contact, modal and buckling, thermal and thermal stress. CFD later\\
FDM materials & Isotropic + knockdown by default; orthotropic by print orientation available per study\\
Drawing standard & ASME Y14.5, third-angle projection\\
Drawings & Associative BOM + balloons; GD\&T and datums with model-driven callouts; revision table and
  release; auto-generate drawing\\
PDM home & \textbf{Evaluate first} (Ops as PLM vs Ops + Git LFS vs in-FreeCAD vault vs open-source PLM)\\
Lifecycle & WIP $\rightarrow$ In Review $\rightarrow$ Released $\rightarrow$ Obsolete. Numeric pre-release
  revisions (01, 02, \ldots), letter at first release (A), then ECR $\rightarrow$ ECO for every later letter
  (ASME Y14.35 + SW-PDM practice)\\
Revision vs P/N & Revision is metadata, separate from the part number; drawings and packages print
  \code{P/N-REV}\\
Team & 2--5 concurrent editors: check-out locks, review and approval\\
Outputs & 3D-print prep; CNC CAM; sheet-metal flat DXF; release package\\
Kept idea areas & AI / MCP assistant; KiCad bridge; robotics and drone kits; kernel study\\
UI items & Patterns $\rightarrow$ Modeling; workspaces, panels, toolbar $\rightarrow$ UI track (later);
  command line $\rightarrow$ Automation\\
Upstream & \textbf{Track FreeCAD main} continuously (replaces ``stay on 1.1.4, rebase at 26.3'')\\
Deliverable & This document plus \file{roadmap-epics.csv}; epics only\\
Rule (standing) & Every feature is on by default, with an off switch on Preferences $\rightarrow$ BruceWorx\\
\bottomrule
\end{longtable}
\end{center}

\subsection{The lifecycle and where parity stands}

\begin{figure}[ht]
\centering
\begin{tikzpicture}[>=Stealth,
  st/.style={draw=bwblue!70,fill=bwblue!6,rounded corners=3pt,minimum width=2.15cm,minimum height=1.05cm,
             align=center,font=\small},
  lab/.style={font=\scriptsize\color{bwgrey},align=center}]
\node[st] (s1) {Sketch \&\\part};
\node[st,right=0.35cm of s1] (s2) {Assembly};
\node[st,right=0.35cm of s2] (s3) {Simulate};
\node[st,right=0.35cm of s3] (s4) {Drawings};
\node[st,right=0.35cm of s4] (s5) {Manage \&\\release};
\node[st,right=0.35cm of s5] (s6) {Make};
\foreach \a/\b in {s1/s2,s2/s3,s3/s4,s4/s5,s5/s6} \draw[->,thick,bwblue!70] (\a) -- (\b);
\draw[->,bwgrey,dashed] (s6.south) -- ++(0,-0.45) -| node[pos=0.25,below,lab]{ECR $\rightarrow$ ECO: the next revision re-enters the loop} (s1.south);
\node[lab,above=0.12cm of s1] {MOD};
\node[lab,above=0.12cm of s2] {ASM, ROB, ECD};
\node[lab,above=0.12cm of s3] {SIM};
\node[lab,above=0.12cm of s4] {DRW};
\node[lab,above=0.12cm of s5] {PDM};
\node[lab,above=0.12cm of s6] {MFG};
\node[lab,below=0.95cm of s1] {\phalf};
\node[lab,below=0.95cm of s2] {\phalf};
\node[lab,below=0.95cm of s3] {\phalf};
\node[lab,below=0.95cm of s4] {\phalf};
\node[lab,below=0.95cm of s5] {\pnone};
\node[lab,below=0.95cm of s6] {\phalf};
\end{tikzpicture}
\caption{Lifecycle stages, the tracks that serve them, and FreeCAD~1.1.4's overall standing against
Fusion/SolidWorks (\pfull{} full, \phalf{} partial, \pnone{} none). Part and drawing management is the only
stage with nothing in core FreeCAD. Foundation (FND), automation (AUT), UI (UIX) and kernel (KRN) cut across
all stages.}
\label{fig:lifecycle}
\end{figure}

The detailed matrix (Appendix~\ref{app:parity}, 71 capabilities) shows the pattern behind
Figure~\ref{fig:lifecycle}:
\begin{itemize}
  \item \textbf{Strong already:} parametric solids, joints and kinematic simulation (Assembly in 1.1),
        CalculiX/Elmer analysis types, drawing views, CAM and post-processors, 3MF and DXF export, Python.
  \item \textbf{Partial:} patterns (with a bug, \S\ref{sec:bug}), parameters, direct edit (Defeaturing only),
        surfacing, threads (drawing-side only), BOM (four automatic columns), GD\&T (symbol library, not
        semantic), mesh $\rightarrow$ solid (faceted).
  \item \textbf{Missing:} configurations, freeform, weldments, interference, smart fasteners, dynamic
        motion, safety-factor and orthotropic results, auto drawings, \emph{all} of PDM, slicer handoff,
        rendering, generative design.
\end{itemize}

\subsection{Milestones}

Milestones are \emph{gates}, not calendar dates. Tracks run in parallel. A milestone closes when all of its
epics pass acceptance on real BruceWorx parts. The Leonardo pedestal assembly \code{LE-02001} is the pilot
throughout.

\begin{center}
\small
\begin{tabularx}{\linewidth}{@{}lp{4.3cm}X@{}}
\toprule
& Milestone & Exit criterion\\
\midrule
M0 & Foundation green & CI green on Linux and macOS arm64; tagged release; nightly rebase on upstream main; review fixes in\\
M1 & Model it & LE-02001 modelled in FreeCAD from scratch: one file per part, patterns, configurations, threads, joints, interference-free\\
M2 & Document \& release it & LE-02001 rev A released \emph{from FreeCAD}: ASME drawings, BOM, PDM lifecycle, release package\\
M3 & Make it & Printed, machined and laser-cut parts made from FreeCAD outputs without touching Fusion\\
M4 & Validate it & Static, modal, buckling and thermal studies on LE parts, verified against benchmarks\\
M5 & Advanced & Surfacing, freeform, mesh $\rightarrow$ solid, in-context, large-assembly, dynamic motion, command line, assistant\\
M6 & Switch system of record & Parity matrix met for the LE scope; Fusion data migrated; FC-T01 criteria satisfied\\
UI & UI and styling pass & After M2: workspaces, panels, fresh visual design, regenerated figures\\
BG & Background & Kernel study; runs at low priority throughout\\
\bottomrule
\end{tabularx}
\end{center}

\begin{figure}[ht]
\centering
% Generated by scripts/epics.py; do not edit.
\begin{tabular}{@{}>{\raggedright\arraybackslash}p{2.35cm}*{9}{>{\centering\arraybackslash}p{1.12cm}}@{}}
\toprule
\textbf{Track} & \textbf{M0} & \textbf{M1} & \textbf{M2} & \textbf{M3} & \textbf{M4} & \textbf{M5} & \textbf{M6} & \textbf{UI} & \textbf{BG}\\
\midrule
\scriptsize Foundation \& upstream & \scriptsize \textcolor{bworange}{01} \textcolor{bworange}{02} \textcolor{bwdark}{03} \textcolor{bwdark}{04} \textcolor{bwdark}{05} \textcolor{bworange}{06} \textcolor{bwgreen}{07} &  &  &  &  &  &  &  & \\[2pt]
\scriptsize Part modeling &  & \scriptsize \textcolor{bworange}{01} \textcolor{bwdark}{02} \textcolor{bwdark}{03} \textcolor{bwdark}{04} \textcolor{bwdark}{05} \textcolor{bwdark}{11} \textcolor{bwdark}{12} &  & \scriptsize \textcolor{bwdark}{06} \textcolor{bwdark}{07} &  & \scriptsize \textcolor{bwdark}{08} \textcolor{bwdark}{09} \textcolor{bwdark}{10} &  &  & \\[2pt]
\scriptsize Assembly &  & \scriptsize \textcolor{bwdark}{01} \textcolor{bwgreen}{02} \textcolor{bworange}{03} \textcolor{bworange}{04} \textcolor{bwdark}{05} \textcolor{bwdark}{06} &  &  &  & \scriptsize \textcolor{bwdark}{07} \textcolor{bwdark}{08} \textcolor{bwdark}{09} \textcolor{bwdark}{10} &  &  & \\[2pt]
\scriptsize Drawings (ASME) &  &  & \scriptsize \textcolor{bwdark}{01} \textcolor{bwdark}{02} \textcolor{bwdark}{03} \textcolor{bwdark}{04} \textcolor{bwdark}{05} \textcolor{bwdark}{06} \textcolor{bwdark}{08} & \scriptsize \textcolor{bwdark}{07} &  &  &  &  & \\[2pt]
\scriptsize Part \& drawing management &  &  & \scriptsize \textcolor{bwgreen}{01} \textcolor{bworange}{02} \textcolor{bworange}{03} \textcolor{bwdark}{04} \textcolor{bwdark}{05} \textcolor{bwdark}{06} \textcolor{bwdark}{07} &  &  &  & \scriptsize \textcolor{bwdark}{08} &  & \\[2pt]
\scriptsize Simulation &  &  &  &  & \scriptsize \textcolor{bwdark}{01} \textcolor{bwdark}{02} \textcolor{bwdark}{03} \textcolor{bwdark}{04} \textcolor{bwdark}{05} & \scriptsize \textcolor{bwdark}{06} &  &  & \\[2pt]
\scriptsize Manufacturing outputs &  &  & \scriptsize \textcolor{bwdark}{04} & \scriptsize \textcolor{bwdark}{01} \textcolor{bwdark}{02} \textcolor{bwdark}{03} &  &  &  &  & \\[2pt]
\scriptsize Robotics \& drone kits &  & \scriptsize \textcolor{bwdark}{01} \textcolor{bwdark}{03} &  & \scriptsize \textcolor{bwdark}{02} \textcolor{bwdark}{04} &  & \scriptsize \textcolor{bwdark}{05} &  &  & \\[2pt]
\scriptsize ECAD bridge (KiCad) &  & \scriptsize \textcolor{bwdark}{01} &  & \scriptsize \textcolor{bwdark}{02} \textcolor{bwdark}{03} \textcolor{bwdark}{04} &  &  &  &  & \\[2pt]
\scriptsize Wire routing &  &  &  & \scriptsize \textcolor{bwdark}{01} \textcolor{bwdark}{02} \textcolor{bwdark}{03} \textcolor{bwdark}{04} &  & \scriptsize \textcolor{bwdark}{05} \textcolor{bwdark}{06} &  &  & \\[2pt]
\scriptsize AI \& automation & \scriptsize \textcolor{bwgreen}{01} & \scriptsize \textcolor{bwdark}{02} \textcolor{bwdark}{03} &  &  &  & \scriptsize \textcolor{bwdark}{04} \textcolor{bwdark}{05} &  &  & \\[2pt]
\scriptsize UI \& styling &  &  &  &  &  &  &  & \scriptsize \textcolor{bwgreen}{01} \textcolor{bwgreen}{02} \textcolor{bworange}{03} \textcolor{bworange}{04} \textcolor{bwdark}{05} \textcolor{bwdark}{06} \textcolor{bwdark}{07} \textcolor{bwdark}{08} \textcolor{bworange}{09} & \\[2pt]
\scriptsize Kernel study &  &  &  &  &  &  &  &  & \scriptsize \textcolor{bwdark}{01} \textcolor{bwdark}{02} \textcolor{bwdark}{03}\\[2pt]
\bottomrule
\end{tabular}

\caption{Epics by track and milestone (number = epic, e.g.\ Part modeling \code{03} = \eref{BW-MOD-03}).
Colour shows status: \textcolor{bwgreen}{\textbf{done}}, \textcolor{bworange}{\textbf{in progress}}, to do.
Generated from \file{scripts/epics.py}.}
\label{fig:swim}
\end{figure}

\subsection{Critical path}

\begin{figure}[ht]
\centering
\begin{tikzpicture}[>=Stealth,x=2.7cm,y=-0.9cm,
  e/.style={draw=black!30,rounded corners=2pt,fill=white,font=\scriptsize\ttfamily,minimum width=1.5cm,minimum height=0.55cm},
  m/.style={draw=bwblue,fill=bwblue!10,rounded corners=2pt,font=\scriptsize\bfseries,minimum width=0.9cm,minimum height=0.55cm}]
\node[e] (f6) at (0,0) {FND-06};  \node[e] (m1) at (1,0) {MOD-01};  \node[m] (M1) at (3,0) {M1};
\node[e] (a1) at (0,1) {ASM-01};  \node[e] (a4) at (1,1) {ASM-04};  \node[e] (d2) at (2,1) {DRW-02};
\node[e] (p6) at (2,2) {PDM-06};  \node[e] (o4) at (3,2) {MFG-04};
\node[e] (p1) at (0,3) {PDM-01};  \node[e] (p4) at (1,3) {PDM-04};  \node[e] (p5) at (2,3) {PDM-05};
\node[e] (p7) at (3,3) {PDM-07};  \node[m] (M2) at (4,3) {M2};
\node[e] (p2) at (0,4) {PDM-02};  \node[e] (d5) at (2,4) {DRW-05};
\node[e] (p3) at (1,5) {PDM-03};
\node[e] (m5) at (0,6) {MOD-05};  \node[e] (d4) at (1,6) {DRW-04};  \node[e] (d6) at (2,6) {DRW-06};
\draw[->] (f6) -- (m1); \draw[->] (m1) -- (M1); \draw[->] (a4) -- (M1);
\draw[->] (a1) -- (a4); \draw[->] (a4) -- (d2); \draw[->] (a4) -- (p6); \draw[->] (p4) -- (p6);
\draw[->] (p1) -- (p4); \draw[->] (p2) -- (p4); \draw[->] (p2) -- (p3); \draw[->] (p2) -- (d5);
\draw[->] (p4) -- (p5); \draw[->] (p5) -- (p7); \draw[->] (d5) -- (p7); \draw[->] (p5) -- (o4);
\draw[->] (p7) -- (M2); \draw[->] (o4) -- (M2);
\draw[->] (m5) -- (d4); \draw[->] (d4) -- (d6);
\end{tikzpicture}
\caption{The dependencies that decide when FreeCAD can release a documented assembly (M2). The PDM
evaluation (\eref{PDM-01}) and data model (\eref{PDM-02}) have no prerequisites, which is why they start on
day one alongside modeling.}
\label{fig:critical}
\end{figure}

\subsection{All epics}
\small
% Generated by scripts/epics.py; do not edit.
\begin{longtable}{@{}p{1.4cm}p{5.8cm}p{0.9cm}p{4.3cm}p{1.95cm}@{}}
\toprule ID & Epic & Mile. & Depends on & Status\\ \midrule\endhead
\bottomrule\endfoot
\midrule\multicolumn{5}{@{}l}{\textbf{Foundation \& upstream}}\\
\code{FND-01} & Build \& CI on Linux and macOS arm64 & M0 & --- & \stProg\\
\code{FND-02} & Packaging, branding \& releases & M0 & \scriptsize FND-01 & \stProg\\
\code{FND-03} & Track upstream FreeCAD main & M0 & \scriptsize FND-01 & \stTodo\\
\code{FND-04} & Licences \& notices & M0 & --- & \stTodo\\
\code{FND-05} & Acceptance corpus \& test harness & M0 & --- & \stTodo\\
\code{FND-06} & Fixes found in review & M0 & --- & \stProg\\
\code{FND-07} & On by default, each with an off switch & M0 & --- & \stDone\\
\midrule\multicolumn{5}{@{}l}{\textbf{Part modeling}}\\
\code{MOD-01} & Rectangular \& circular patterns & M1 & \scriptsize FND-06 & \stProg\\
\code{MOD-02} & Parameters \& top-down design & M1 & --- & \stTodo\\
\code{MOD-03} & Configurations & M1 & \scriptsize MOD-02 & \stTodo\\
\code{MOD-04} & Direct edit on imported geometry & M1 & --- & \stTodo\\
\code{MOD-05} & Threads, holes \& hardware library & M1 & --- & \stTodo\\
\code{MOD-06} & Sheet metal & M3 & --- & \stTodo\\
\code{MOD-07} & Frames \& weldments & M3 & --- & \stTodo\\
\code{MOD-08} & Surfacing parity & M5 & --- & \stTodo\\
\code{MOD-09} & Freeform (subdivision / T-spline-like) & M5 & \scriptsize MOD-08 & \stTodo\\
\code{MOD-10} & Mesh to solid & M5 & --- & \stTodo\\
\code{MOD-11} & Inspect tools & M1 & --- & \stTodo\\
\code{MOD-12} & Sketcher parity audit & M1 & --- & \stTodo\\
\midrule\multicolumn{5}{@{}l}{\textbf{Assembly}}\\
\code{ASM-01} & File model: one part per file & M1 & --- & \stTodo\\
\code{ASM-02} & Assembly gap audit & M1 & --- & \stDone\\
\code{ASM-03} & Joint UX \& limits & M1 & \scriptsize ASM-02 & \stProg\\
\code{ASM-04} & Exploded views \& BOM & M1 & \scriptsize ASM-01 & \stProg\\
\code{ASM-05} & Sub-assemblies & M1 & \scriptsize ASM-01 & \stTodo\\
\code{ASM-06} & Interference \& clearance & M1 & --- & \stTodo\\
\code{ASM-07} & In-context editing & M5 & \scriptsize ASM-01 MOD-02 & \stTodo\\
\code{ASM-08} & Large-assembly mode & M5 & \scriptsize ASM-01 & \stTodo\\
\code{ASM-09} & Smart fastener mating & M5 & \scriptsize MOD-05 ASM-03 & \stTodo\\
\code{ASM-10} & Dynamic motion study & M5 & \scriptsize ASM-03 & \stTodo\\
\midrule\multicolumn{5}{@{}l}{\textbf{Drawings (ASME)}}\\
\code{DRW-01} & ASME templates \& title block & M2 & --- & \stTodo\\
\code{DRW-02} & Associative BOM table \& balloons & M2 & \scriptsize ASM-04 & \stTodo\\
\code{DRW-03} & GD\&T \& datums (ASME Y14.5) & M2 & --- & \stTodo\\
\code{DRW-04} & Hole \& thread callouts from the model & M2 & \scriptsize MOD-05 & \stTodo\\
\code{DRW-05} & Revision table \& release stamp & M2 & \scriptsize PDM-02 & \stTodo\\
\code{DRW-06} & Auto-generate drawing & M2 & \scriptsize DRW-01 DRW-04 & \stTodo\\
\code{DRW-07} & Sheet-metal flat \& cut-list drawings & M3 & \scriptsize MOD-06 MOD-07 & \stTodo\\
\code{DRW-08} & Publish PDF/DXF & M2 & \scriptsize DRW-01 & \stTodo\\
\midrule\multicolumn{5}{@{}l}{\textbf{Part \& drawing management}}\\
\code{PDM-01} & PDM evaluation \& decision & M2 & --- & \stDone\\
\code{PDM-02} & Data model \& lifecycle spec & M2 & --- & \stProg\\
\code{PDM-03} & Numbering service & M2 & \scriptsize PDM-02 & \stProg\\
\code{PDM-04} & FreeCAD vault client & M2 & \scriptsize PDM-01 PDM-02 & \stTodo\\
\code{PDM-05} & Lifecycle \& change control & M2 & \scriptsize PDM-04 & \stTodo\\
\code{PDM-06} & BOM sync \& where-used & M2 & \scriptsize ASM-04 PDM-04 & \stTodo\\
\code{PDM-07} & Pilot: LE-02001 \& CAD register & M2 & \scriptsize PDM-05 DRW-05 & \stTodo\\
\code{PDM-08} & Migration from Fusion & M6 & \scriptsize PDM-07 & \stTodo\\
\midrule\multicolumn{5}{@{}l}{\textbf{Simulation}}\\
\code{SIM-01} & Study setup (Fusion-style) & M4 & --- & \stTodo\\
\code{SIM-02} & Contacts, bolts \& preload & M4 & \scriptsize SIM-01 & \stTodo\\
\code{SIM-03} & Materials incl. FDM & M4 & \scriptsize SIM-01 & \stTodo\\
\code{SIM-04} & Results \& reports & M4 & \scriptsize SIM-01 & \stTodo\\
\code{SIM-05} & Verification benchmarks & M4 & \scriptsize SIM-01 & \stTodo\\
\code{SIM-06} & CFD / airflow (later) & M5 & \scriptsize SIM-01 & \stTodo\\
\midrule\multicolumn{5}{@{}l}{\textbf{Manufacturing outputs}}\\
\code{MFG-01} & 3D-print prep & M3 & \scriptsize MOD-11 & \stTodo\\
\code{MFG-02} & CNC CAM & M3 & --- & \stTodo\\
\code{MFG-03} & Sheet-metal flat DXF & M3 & \scriptsize MOD-06 & \stTodo\\
\code{MFG-04} & Release package & M2 & \scriptsize PDM-05 DRW-08 & \stTodo\\
\midrule\multicolumn{5}{@{}l}{\textbf{Robotics \& drone kits}}\\
\code{ROB-01} & COTS library \& insert & M1 & \scriptsize MOD-05 & \stTodo\\
\code{ROB-02} & Frame \& arm templates & M3 & \scriptsize MOD-03 & \stTodo\\
\code{ROB-03} & Mass, CG \& inertia report & M1 & \scriptsize MOD-11 ROB-01 & \stTodo\\
\code{ROB-04} & Drone balance \& thrust-to-weight & M3 & \scriptsize ROB-03 & \stTodo\\
\code{ROB-05} & Kinematic chains \& joint drive & M5 & \scriptsize ASM-03 & \stTodo\\
\midrule\multicolumn{5}{@{}l}{\textbf{ECAD bridge (KiCad)}}\\
\code{ECD-01} & ECAD intake \& adopt-vs-build & M1 & --- & \stTodo\\
\code{ECD-02} & Board STEP into assembly & M3 & \scriptsize ECD-01 & \stTodo\\
\code{ECD-03} & Board outline round-trip & M3 & \scriptsize ECD-01 & \stTodo\\
\code{ECD-04} & Placement sync \& enclosure fit & M3 & \scriptsize ECD-02 ASM-06 & \stTodo\\
\midrule\multicolumn{5}{@{}l}{\textbf{Wire routing}}\\
\code{WIR-01} & Adopt versus build & M3 & --- & \stTodo\\
\code{WIR-02} & Port and wire & M3 & \scriptsize WIR-01 & \stTodo\\
\code{WIR-03} & Assembly follow & M3 & \scriptsize WIR-02 ASM-03 & \stTodo\\
\code{WIR-04} & From-to cut list & M3 & \scriptsize WIR-02 ASM-04 & \stTodo\\
\code{WIR-05} & Clips and bundles & M5 & \scriptsize WIR-03 & \stTodo\\
\code{WIR-06} & Route on the clip graph & M5 & \scriptsize WIR-05 & \stTodo\\
\midrule\multicolumn{5}{@{}l}{\textbf{AI \& automation}}\\
\code{AUT-01} & Headless hooks \& CLI & M0 & --- & \stDone\\
\code{AUT-02} & bruceworx.api & M1 & \scriptsize AUT-01 & \stTodo\\
\code{AUT-03} & MCP server & M1 & \scriptsize AUT-02 & \stTodo\\
\code{AUT-04} & Assistant panel & M5 & \scriptsize AUT-03 & \stTodo\\
\code{AUT-05} & Command line & M5 & --- & \stTodo\\
\midrule\multicolumn{5}{@{}l}{\textbf{UI \& styling}}\\
\code{UIX-01} & Theme, pref pack \& viewport styling & UI & --- & \stDone\\
\code{UIX-02} & macOS title bar & UI & --- & \stDone\\
\code{UIX-03} & Browser (overlay model tree) & UI & --- & \stProg\\
\code{UIX-04} & History timeline & UI & --- & \stProg\\
\code{UIX-05} & Workspaces \& tabs (Fusion naming) & UI & --- & \stTodo\\
\code{UIX-06} & Toolbar \& panels & UI & \scriptsize UIX-05 & \stTodo\\
\code{UIX-07} & Task panel polish & UI & --- & \stTodo\\
\code{UIX-08} & Matcap shading & UI & --- & \stTodo\\
\code{UIX-09} & Fresh UI design \& re-render figures & UI & \scriptsize UIX-05 & \stProg\\
\midrule\multicolumn{5}{@{}l}{\textbf{Kernel study}}\\
\code{KRN-01} & OCCT coupling survey & BG & --- & \stTodo\\
\code{KRN-02} & STEP exchange harness & BG & \scriptsize FND-05 & \stTodo\\
\code{KRN-03} & Kernel decision & BG & \scriptsize KRN-01 KRN-02 & \stTodo\\
\end{longtable}

\normalsize
